\documentclass[lettersize,journal]{IEEEtran}
\usepackage{amsmath,amsfonts}
% \usepackage{algorithmic}
\usepackage{algpseudocode}
% \usepackage{algorithm}
\usepackage{array}
\usepackage[caption=false,font=normalsize,labelfont=sf,textfont=sf]{subfig}
\usepackage{textcomp}
\usepackage{stfloats}
\usepackage{url}
\usepackage{verbatim}
\usepackage{graphicx}
\usepackage{cite}
\hyphenation{op-tical net-works semi-conduc-tor IEEE-Xplore}
% updated with editorial comments 8/9/2021
\usepackage[hidelinks]{hyperref}%文档插入超链接
\hypersetup{
    colorlinks=false,   %set true if you want colored links
    linktoc=all,        %set to all if you want both sections and subsections linked
    linkcolor=black,     %choose some color if you want links to stand out
}
\usepackage{booktabs}
\usepackage{wrapfig}
% \usepackage[algochapter,linesnumbered,ruled,lined,boxed]{algorithm2e}
\usepackage[linesnumbered,ruled,lined,boxed]{algorithm2e}
\usepackage[numbers,sort&compress]{natbib}
% \biboptions{sort&compress}
\usepackage{url}


\begin{document}

% \title{A Contrastive Learning Method Based on Hard Sample Mining for Wind Turbine Pitch System Fault Diagnosis}
% \title{A Fault Diagnosis Method of Wind Turbine Pitch System Combining Contrastive Learning and Hard Sample Mining}
% \title{A Supervised Contrastive Learning Method for Wind Turbine Pitch System Fault Diagnosis Enabled by Hard Sample Mining}
\title{Hard Sample Mining Enabled Supervised Contrastive Feature Learning for Wind Turbine Pitch System Fault Diagnosis}

% \author{Zixuan Wang\fnref{fn1}}
% \ead{zixuanw.20@intl.zju.edu.cn}
% \author{Bo Qin\fnref{fn1}}
% \ead{bo.20@intl.zju.edu.cn}
% \author{Peng Peng\corref{mycorrespondingauthor}}
% \ead{pengp17@mails.tsinghua.edu.cn}
% \author{Hongwei Wang\corref{mycorrespondingauthor}}
% \ead{hongweiwang@zju.edu.cn}

\author{Zixuan Wang, Bo Qin, Mengxuan Li, Chenlu Zhan, Mark D.\ Butala, Peng Peng,~\IEEEmembership{Member,~IEEE}, and Hongwei Wang,~\IEEEmembership{Member,~IEEE}
        % <-this % stops a space
\thanks{Zixuan Wang and Bo Qin contribute equally to this work.}
\thanks{Zixuan Wang is with the College of Biomedical Engineering and Instrument Science in Zhejiang University, Hangzhou, 310013, China. (E-mail: \href{mailto:zixuanw.20@intl.zju.edu.cn}{zixuanw.20@intl.zju.edu.cn}.}
\thanks{Mengxuan Li and Chenlu Zhan are with the College of Computer Science and Technology in Zhejiang University, Hangzhou, 310013, China. (E-mail: \href{mailto:mengxuanli@intl.zju.edu.cn}{mengxuanli@intl.zju.edu.cn}).}
\thanks{Bo Qin, Mark D.\ Butala, Peng Peng and Hongwei Wang are with Zhejiang University and the University of Illinois Urbana–Champaign Institute, Haining, 314400, China. (E-mail:  \href{mailto:bo.20@intl.zju.edu.cn}{bo.20@intl.zju.edu.cn}, 
\href{mailto:markbutala@intl.zju.edu.cn}{markbutala@intl.zju.edu.cn}, \href{mailto:pengpeng@intl.zju.edu.cn}{pengpeng@intl.zju.edu.cn}, \href{mailto:hongweiwang@intl.zju.edu.cn}{hongweiwang@intl.zju.edu.cn}).}
% \thanks{This work was supported by the National Key RD Program of China (No.2020YFB1707803) and National Natural Science Foundation of China (No.62276230). (Corresponding authors: Peng Peng and Hongwei Wang)}
\thanks{Corresponding authors: Peng Peng and Hongwei Wang}
% \thanks{This work has been submitted to the IEEE Transactions on Sustainable Energy for possible publication. Copyright may be transferred without notice, after which this version may no longer be accessible.}% <-this % stops a space
        }

% The paper headers
% \markboth{Journal of \LaTeX\ Class Files,~Vol.~14, No.~8, August~2021}%
% {Shell \MakeLowercase{\textit{et al.}}: A Sample Article Using IEEEtran.cls for IEEE Journals}

% \IEEEpubid{0000--0000/00\$00.00~\copyright~2021 IEEE}
% Remember, if you use this you must call \IEEEpubidadjcol in the second
% column for its text to clear the IEEEpubid mark.

\maketitle

\begin{abstract}
%原来的版本
% The pitch system of a wind turbine can adjust the blade pitch angle to ensure that the wind turbine can efficiently utilize wind power at different wind speeds. However, there are multiple types of faults in the pitch system that can affect the operation of the wind turbine, and these faults are difficult to classify among themselves. In this paper, a hard sample mining-enabled contrastive feature learning method is proposed to solve this problem. The proposed method first determines the hard samples by the cosine similarity between samples and then utilizes contrastive feature learning to learn a better representation by constructing hard sample pairs. Subsequently, a multilayer perceptron is trained based on the learned discriminative representations to be an efficient classifier. The effectiveness of the proposed method is validated on two real datasets of wind turbine pitch system cog belt fracture datasets, and its fault diagnosis performance far exceeds that of the compared methods.
%chatgpt润色的版本
The efficient utilization of wind power by wind turbines relies on the ability of their pitch systems to adjust blade pitch angles in response to varying wind speeds. However, the presence of multiple health conditions in the pitch system due to the long-term wear and tear poses challenges in accurately classifying them, thus increasing the maintenance cost of wind turbines or even damaging them. This paper proposes a novel method based on hard sample mining-enabled supervised contrastive learning (HSMSCL) to address this problem. The proposed method employs cosine similarity to identify hard samples and subsequently, leverages supervised contrastive learning to learn more discriminative representations by constructing hard sample pairs. 
\begin{comment}
    Furthermore, a multilayer perceptron (MLP) is trained using the learned representations to serve as an efficient classifier.
\end{comment}
Furthermore, the hard sample mining framework in the proposed method also constructs hard samples with learned representations to make the training process of the multilayer perceptron (MLP) more challenging and make it a more effective classifier.
The proposed approach progressively improves the fault diagnosis model by introducing hard samples in the SCL and MLP phases, thus enhancing its performance in complex multi-class fault diagnosis tasks.

To evaluate the effectiveness of the proposed method, two real datasets comprising wind turbine pitch system cog belt fracture data are utilized. The fault diagnosis performance of the proposed method is compared against existing methods, and the results demonstrate its superior performance. The proposed approach exhibits significant improvements in fault diagnosis performance, providing promising prospects for enhancing the reliability and efficiency of wind turbine pitch system fault diagnosis.

\end{abstract}

\begin{IEEEkeywords}
Wind turbine, Pitch system, Fault diagnosis, Hard sample mining, Class imbalance
\end{IEEEkeywords}

\section{Introduction}%介绍风机变桨系统
\IEEEPARstart{W}{ind} energy has been widely used for decades. According to the Global Wind Energy Council (GWEC) statistic in 2022, nearly 94 GW of new wind turbine capacity was installed in 2021, bringing the global wind power capacity to 837 GW \cite{gwec2022global}. The maintenance cost is also increasing with the increasing number of wind turbines. Therefore, condition monitoring and fault diagnosis for wind turbines is crucial to control the operating costs of wind turbines. 

A wind turbine is a complex electromechanical system consisting of blades, rotors, pitch system, yaw system, gearbox, generator, and other components \cite{qiao2015survey}. According to a survey conducted by the China Wind Energy Association (CWEA) between 2010 and 2012 among 47 wind turbine manufacturers, component suppliers and developers, damage to these components is the main cause of wind turbine failures. Among them, the frequency of failure of the pitch system ranked high \cite{lin2016fault}. 

The general structure of the wind turbine pitch system is shown in Fig.\ref{pitchstructure}. The Pitch system can adjust the blade pitch angle in real-time according to the change of wind speed, so as to make full use of wind energy for power generation and reduce the impact of wind on the blade \cite{perez2013wind}. In addition, the pitch system can feather the blade to achieve a safe shutdown when the wind speed is too high or during an emergency shutdown \cite{hansen2004review}. There are two types of pitch systems for wind turbines: electric motor drive and hydraulic drive. Due to the oil leakage and maintenance problems of hydraulic drive, the pitch systems of wind turbines in China are mainly driven by electric motors \cite {lin2016fault}. The main failures of electric motor-driven pitch systems include electric motor failure, transmission component failure, and control system failure. Therefore, the timely diagnosis of wind turbine pitch system faults can effectively reduce downtime and losses. 

\begin{figure}[!t]
\centering
\includegraphics[width=0.95\linewidth]{Figure/methodology/wind_turbine_pitch_system_structure_new.pdf}
\caption{The general structure of the wind turbine pitch system}
\label{pitchstructure}
\end{figure}

%此部分考虑介绍风机故障诊断分为三类：信号处理的方法、基于模型的方法以及数据驱动的方法。related works中介绍风机变桨系统故障诊断的数据驱动方法，如此便可不重复。
% \textbf{The fault diagnosis methods of wind turbine pitch systems are mainly divided into model-based methods and data-driven methods. (not sure about the classes).} 

The methods of wind turbine fault diagnosis are mainly divided into signal analysis-based methods, model-based methods, and data-driven methods \cite{helbing2018deep,li2021wind,habibi2019reliability}.
Signal analysis-based methods acquire data through sensors installed in various components of the wind turbine and extract fault-related features for fault diagnosis \cite{wang2019vibration,liu2020review,salameh2018gearbox}.  For example, \citet{an2012application} put forward a fault diagnosis method of wind turbine bearing based on intrinsic time-scale decomposition. \citet{zhang2018wavelet} proposed a wavelet energy transmissibility function that is robust to noise and sensitive to system property changes to diagnose the fault in wind turbine bearings. \citet{tang2010wind} proposed an effective wind turbine fault diagnosis method with good denoising performance based on the Morlet wavelet transformation and Wigner-Ville distribution. Although signal analysis-based methods possess high fault detection accuracy, they are extremely dependent on additional sensors, which themselves will inevitably fail, thus affecting the accurate diagnosis of faults. In addition, the arrangement of sensors also increases the operational cost and wiring complexity of the wind turbine \cite{qiao2015survey2}. 

Model-based approaches construct a quantitative or qualitative model for fault diagnosis through prior knowledge of the operating modes of the wind turbine system \cite{abichou2014fault}. For example, \citet{Chaaban2014} modeled a utility-scale turbine using an aero-hydro-servo-elastic simulation tool to simulate the faults inside the pitch system and to study their effects as a function of the fault magnitude and wind speed. \citet{odgaard2012model} proposed a fault-tolerant observer scheme to provide estimates of various components for a wind turbine. \citet{lan2018fault} proposed an adaptive sliding observer to estimate wind turbine parametric pitch actuator faults. Model-based wind turbine fault diagnosis techniques do not rely on additional hardware components, but their diagnostic performance depends entirely on the mathematical model. The complex structure of a wind turbine and its dynamic operating states make it particularly challenging to construct a suitable mathematical model \cite{9774410}. 

Data-driven methods use a large amount of historical data which can be obtained through the supervisory control and data acquisition (SCADA) system installed in the wind turbine to train the fault diagnosis model. For example, \citet{gao2018novel} proposed a wind turbine fault diagnosis method based on integral extension load mean decomposition multiscale entropy and least squares support vector machine.  \citet{lei2019fault} proposed an end-to-end Long Short-term Memory (LSTM) model for fault diagnosis of wind turbines. \citet{pashazadeh2018data} proposed a data-driven approach based on the fusion of several classifiers for fault detection and isolation in wind turbines. However, data-driven approaches still face some challenges. (1). The complex and variable operating states of wind turbines lead to a huge amount of data collected by SCADA systems with the characteristics of class imbalance, nonlinearity, and nonstationarity, which makes it difficult to effectively extract features from the data; (2). Existing studies rarely focus on the problem of hard samples in multiple health states which is difficult to distinguish in wind turbine fault diagnosis tasks.
% However, existing data-driven approaches rarely focus on the problem of hard samples which is difficult to distinguish in wind turbine fault diagnosis tasks.

% For example, \citet{pashazadeh2018data} proposed a data-driven approach based on the fusion of several classifiers for fault detection and isolation in wind turbines.\textbf{change another way using SCADA system}. \citet{gao2018novel} proposed a wind turbine fault diagnosis method based on integral extension load mean decomposition multiscale entropy and least squares support vector machine.  \citet{lei2019fault} proposed an end-to-end Long Short-term Memory (LSTM) model for fault diagnosis of wind turbine. 

%之前的版本
% This paper proposes a novel fault diagnosis method for wind turbine pitch systems that combines supervised contrastive learning and hard sample mining. The purpose of the proposed method is to better distinguish the faults in the wind turbine pitch system whose severity varies with time. In order to achieve this purpose, we face two challenges. First, the amount of SCADA data on wind turbines is extremely large, so how to extract discriminative features from SCADA data is the problem we face. Second, the wind turbine pitch system may have multiple faults that evolve over time, and it is difficult to classify the overlap between different fault types. These two challenges are the key issues in building a fault diagnosis model for wind turbine pitch systems based on SCADA data.

%20230724版本
This paper proposes a novel fault diagnosis method for wind turbine pitch systems that combines supervised contrastive learning and hard sample mining. The purpose of the proposed method is to better distinguish the faults in the wind turbine pitch system whose severity varies with time. In order to achieve this purpose, we face two challenges. First, the amount of SCADA data on wind turbines is extremely large, so how to extract discriminative features from SCADA data is the problem we face. Second, there are many different health conditions in the wind turbine pitch system, which may be difficult to distinguish between them due to the high degree of similarity, so how to accurately distinguish between multiple health conditions is another challenge. These two challenges are the key issues in building a fault diagnosis model for wind turbine pitch systems based on SCADA data.

The main contributions of this paper can be summarized as follows.
\begin{itemize}
    \item A two-step approach consisting of supervised contrastive learning and a multi-layer perceptron is used to diagnose faults in wind turbine pitch systems. Supervised contrastive feature learning extracts discriminative features from large amounts of SCADA data. The extracted features are then utilized to classify the data and discriminate faults by the multi-layer perceptron.
    \item A hard sample mining framework based on cosine similarity is used to distinguish the hard samples in wind turbine SCADA data with multiple health conditions. To the best of our knowledge, few similar studies have been conducted in the field of wind turbine pitch system fault diagnosis, and the hard sample mining framework proposed in this paper fills the gap in this field.
    \begin{comment}
        To the best of our knowledge, this paper is the first to introduce a hard sample mining method in the field of wind turbine pitch system fault diagnosis.
    \end{comment}
    \item The proposed hard sample mining method is a unified framework acting on both supervised contrastive learning and multi-layer perceptron.
\end{itemize}
%is to propose a novel two-step method that can overcome the above challenges. For the first challenge, supervised contrastive learning is used to extract discriminative features from SCADA data. For the second challenge, a hard sample mining method based on cosine similarity is used to distinguish hard-to-classify faults with temporal overlap. The wind turbine pitch system fault diagnosis procedure is divided into two stages. The first stage is supervised contrastive learning enabled by hard sample mining. The second stage is a multi-layer perceptron enabled by hard sample mining to classify the data and discriminate faults. It should be noted that the hard sample mining method is a unified framework acting on both supervised contrastive learning and multi-layer perceptron.

 The rest of this paper is organized as follows. First, the background theory is discussed in Section II. Then, the proposed HSMSCL is described in detail in Section III. The data used for the experiments and the analysis of the experimental results are presented in Section IV. Finally, the conclusions of this paper are given in Section V.
%介绍风机各部件的故障率（谷歌学术关键词Survey of failures in wind power systems），指出变桨系统的重要性。参考Wind Turbine Pitch System Condition Monitoring and Fault Detection Based on Optimized Relevance Vector Machine Regression
%详细介绍变桨系统的功能：主要有两种——电机驱动变桨和液压驱动变桨，并介绍这两种变桨系统可能发生的主要故障。
%对变桨系统的故障诊断方法包括哪些类型？？？（参考文献Deep Learning for fault detection in wind turbines）暂定方案：分为基于模型的和数据驱动的。数据驱动方法部分不详细介绍相关论文，留到background theory里related works部分介绍（数据驱动方法中引出SCADA系统）。
%本文就是利用SCADA系统对风机变桨系统进行基于机器学习（？）/数据驱动的故障诊断。

\section{Background Theory}
% \subsection{Related works}%之前的版本
% %介绍基于深度学习数据驱动的风机变桨系统故障诊断相关工作，引出难样本分类问题。
% %NBM是为了解决正常样本和故障样本分布不平衡问题所以才只用正常样本训练NBM，可以从变桨系统不平衡诊断角度来展开介绍几种方法，然后说这些方法只考虑解决不平衡分布问题，却忽略了难分类样本（NBM无法进行多分类诊断），参考论文《Condition Monitoring of Wind Turbine Generator Based on Transfer Learning and One-Class Classifier》.变桨系统大部分时间都是正常的，故障数量占比少，多种数据驱动的故障诊断方法用于检测变桨系统故障——过采样如SMOTE和NBM。NBM的方法无法应对多种类型的故障，而数据重采样的方法会引入噪声或删去数据中的重要信息。
% % In recent years, data-driven methods have been widely used in the field of wind turbine fault diagnosis. \textbf{change the description}

% In practice, the wind turbine pitch system operates in a healthy state most of the time, resulting in a small number of faults collected by the SCADA system. To cope with this problem, researchers often perform data resampling or construct normal behavior models (NBM). Methods of data resampling include oversampling of the minor class \cite{yang2023fault}, undersampling of the major class \cite{lin2017clustering}, and synthetic minority oversampling technique 
%  (SMOTE) \cite{zhang2018imbalanced,9461331}. 
%  \begin{comment}
%      For example, \citet{9461331} used SMOTE to balance SCADA data and decision trees to predict wind turbine alarm conditions. 
%  \end{comment}
%  However, data resampling methods may introduce undesirable noise or delete important information from the data, thus affecting the accuracy of fault diagnosis. NBM methods train the model using only normal data and use the residuals between the measured and model signals as indicators of a possible fault \cite{tautz2017using}. For example, \citet{wei2019wind} predicted failures for wind turbine electric pitch systems by an optimized relevance vector machine-based NBM method. A limitation of the NBM methods is that it is difficult to distinguish multiple faults. %\textbf{In addition, the above-mentioned methods also ignore multiple faults whose severity varies with time, which are very common in pitch systems. }\textbf{Is the limitation repeated?}
 
%  % In this paper, we propose, for the first time, a supervised contrastive feature learning method enabled by a hard sample mining framework that can well diagnose the hard samples in wind turbine pitch systems.%原来的版本

%  In this paper, we propose, for the first time, a supervised contrastive learning method enabled by a hard sample mining framework for wind turbine pitch system fault diagnosis, which successfully solves the hard sample problem of difficulty in distinguishing between multiple health conditions.%20230723新改的版本

% A highly accurate data-driven classification system for the diagnosis of wind turbine pitch system was presented in \cite{godwin2013classification}, resulting in a pitch fault diagnostic accuracy of 87.05\% along with a 42.12\% reduction in pitch fault alarms.
% For example, \citet{wei2019wind} proposed a new normal behavior modeling method to predict wind turbine electric pitch system failures using SCADA information which is particularly effective for online monitoring applications at reasonable computational complexity. \citet{cho2021fault} proposed an approach that is utilized as a hybrid form based on a Kalman filter and an artificial neural network for the blade pitch system in a spar-type floating wind turbine.

\subsection{Related works}%20230801版本，考虑讲不平衡多分类问题
%实际上，风机变桨系统大部分时间运行在正常状态，而变桨系统的磨损是一个逐渐的过程，在彻底损坏前有多种健康状态，因此变桨系统的故障诊断问题是一个典型的不平衡多分类问题。常见的不平衡多分类问题解决方法主要包括数据层面的方法和算法层面的方法，此外还有NBM方法。
%\textbf{imbalanced multi-class fault diagnosis}
%不平衡故障诊断分为数据层面的方法和算法级别的方法还有NBM（NBM是不是算法级别的？存疑），数据层面的方法原理xxx，列举一些文章，缺点是xxx；算法级别的方法原理是xxx，列举一些文章，缺点是可能需要领域内知识；NBM方法原理是XXX，列举一些文章，该方法不需要先验知识，缺点是难以应对多分类问题
In practice, the wind turbine pitch system operates in a normal state most of the time, while the wear and tear of the pitch system is a gradual process, with multiple health states before complete damage. Therefore, the fault diagnosis problem of the pitch system is a typical imbalanced multi-classification problem. The common methods for solving this problem mainly include data-level methods, algorithm-level methods, and normal behavior model-based(NBM) methods \cite{tanha2020boosting,tautz2017using}. Data-level methods make the data balanced by resampling them, including the undersampling of the major class and the oversampling of the minor class \cite{yang2023fault}. For example, \citet{zhang2018imbalanced} proposed a novel weighted minority oversampling approach for imbalanced data problems in real rotating machinery applications. In \cite{zhang2020machinery}, a generative adversarial network-based method was utilized to generate additional realistic fake samples, thus solving the class imbalance problem in machinery temporal vibration data. The main drawback of the data-level approach is that during the process of data resampling, important information in the data may be deleted or undesirable noise may be introduced, thus affecting the accuracy of fault diagnosis.

%Algorithm-level: Cost-sensitive learning and focal loss，参考《A multi-learner neural network approach to wind turbine fault diagnosis with imbalanced data》中introduction的介绍
The algorithm-level approach improves the bias against the majority class by making modifications to the algorithm. For example, \citet{he2020spatio} proposed a spatio-temporal multiscale neural network and adopted the focal loss as the loss function to address the data imbalance problem. A data-driven framework for rotating machinery diagnosis was proposed in \cite{xu2020imbalanced} and the imbalanced operating condition was effectively diagnosed via cost-sensitive learning. However, algorithm-level approaches often require a deep understanding of the task in order to select the appropriate algorithm.

NBM methods train the model using only normal data and use the residuals between the measured and model signals as indicators of a possible fault \cite{tautz2017using}. For example, \citet{wei2019wind} designed an NBM method using optimized relevance vector machine regression to predict wind turbine electric pitch system failures using SCADA information. In order to detect wind turbine generator faults, normal behavior models were proposed in \cite{bi2017detection} based on performance curves. Although NBM methods are capable of coping with the data imbalance problem, a limitation of the NBM methods is their inability to diagnose multiple classes of faults.

In this paper, we innovatively propose a supervised contrastive learning method enabled by a hard sample mining framework for wind turbine pitch system fault diagnosis, which successfully solves the hard sample problem of difficulty in distinguishing between multiple health conditions.

\subsection{Motivation}%详细介绍难样本存在的原因，参考彭鹏TII那篇论文中的表述，从两个方面解释：1.故障是时间剖分的，且样本之间存在极度不均衡的情况；2.绘制一下SCL方法分类出的混淆矩阵，指出这个数据就是存在难分类的现象。以上是我们工作的动机，所以引出了我们的灵感——通过难样本挖掘的方法解决上述问题（参考和彭师兄交流会议1129版）
The components in the pitch system of a wind turbine gradually wear out during an increasing number of operations, making it challenging to accurately diagnose the health condition of the pitch system components. 
To demonstrate the existence of this phenomenon, we selected a real SCADA dataset of a wind turbine pitch system (This dataset is described in detail in Section \ref{secdatadesc}.) to diagnose its five health conditions by SCL+MLP, the results of which are shown in Fig. \ref{motivationfig}. 
From the figure, it can be seen that there is a serious misclassification problem between the normal class and the four faulty classes. This problem may occur because the health conditions of the pitch system components change gradually over time, and in addition, there is a class imbalance as the different health conditions last for different periods of time. This problem may lead to higher maintenance costs or even damage to the whole wind turbine.

Although data-driven fault diagnosis techniques are widely used for wind turbine pitch systems, to the best of our knowledge, few of them have explored the phenomenon of hard samples between pitch system health conditions. In this paper, we propose a hard sample mining framework based on cosine similarity, which can construct hard sample pairs in a mini-batch to make the training stage more challenging and make the distribution of samples in the feature space more compact and discriminative. Combining the proposed hard sample mining framework with SCL-MLP can well solve the hard sample problem in wind turbine pitch systems. The details of the proposed method will be described in Section \ref{proposedmethod}.


\begin{figure}[!t]%tem=0.5,last_best
\centering
\subfloat[]{\includegraphics[width=1.7in]{Figure/Motivation/Motivation_TSNE_SCL23.png}%
\label{motivation_tsne}}
\hfil
\subfloat[]{\includegraphics[width=1.5in]{Figure/Motivation/Motivation_confusion_matrix_SCL+MLP_Num23.png}%
\label{motivation_confusion}}
\caption{Fault diagnosis result of wind turbine pitch system by SCL+MLP. (a) TSNE embedding visualization (b) Confusion matrix}
\label{motivationfig}
\end{figure}

\subsection{Supervised Contrastive Feature Learning}%介绍SCL的理论与逻辑，20230606改成contrastive feature learning

With its powerful feature extraction capability, contrast learning has been widely used as an efficient feature learning method in natural language processing \cite{gunel2020supervised,inoue2020semi}, computer vision \cite{Wang_2021_CVPR,li2021dual} and other fields. 
\begin{comment}
    Contrastive learning has been widely used in computer vision\cite{Wang_2021_CVPR,li2021dual}, natural language processing\cite{gunel2020supervised,inoue2020semi}, and other fields because of its powerful representation learning capability.
\end{comment} 
Supervised contrastive learning selects an anchor and chooses samples with the same label as the anchor as positives and samples with different labels as negatives. For each anchor, more discriminative representations are learned by pushing the anchor closer to the positives and farther away from the negatives in the embedding space \cite{khosla2020supervised}. The supervised contrastive loss is defined as follows:

\begin{equation}
    \mathcal{L}^\text{sup} = \sum_{i \in I} \frac{-1}{|P(i)|} \sum_{p \in P(i)} log \frac{e^{sim({z_i}, {z_p}) / \tau}}{\sum_{a \in A(i)}
    e^{sim({z_i}, {z_a})/\tau} }
    \label{losssup}
\end{equation}
where $I \in D^l$ is a mini-batch of the labeled data, and $\tau \in \mathcal{R}^{+}$ is a scalar temperature parameter. $P(i)\equiv\left\{p \in A(i): \tilde{\boldsymbol{y}}_p=\tilde{\boldsymbol{y}}_i\right\}$ is the set of indices of all positives of the anchor and $|P(i)|$ is its cardinality. 
\begin{comment}
    Supervised contrastive learning makes full use of the information of labels, so models trained based on supervised contrastive loss can learn more compact representations and therefore, more discriminative features.
\end{comment}
Supervised contrastive learning leverages label information effectively, enabling models trained with supervised contrastive loss to acquire compact representations. Consequently, these models can capture more discriminative features. 
%\textbf{The target function of a model using SCL is:}

\subsection{Hard Sample Mining}%介绍HSM的理论与逻辑

The boundaries of samples of different classes in the feature space may be close to each other, so samples near the boundaries will be difficult to classify, and such samples are hard samples. A schematic diagram of hard samples is shown in Fig. \ref{hardsample}, where the two classes are represented by green and blue dots, respectively, and the data points in the region framed by the red dotted line are the hard samples. Moreover, if a sample is selected as the query sample, the hard samples of the query sample include hard positive samples with the same label but far away from it and hard negative samples with different labels but close to it \cite{smirnov2018hard,chen2020hard}. Making full use of hard samples in the training phase of the model not only improves the learning efficiency but also enables the model to have better classification performance. 

\begin{figure}[!t]
\centering
\includegraphics[width=0.95\linewidth]{Figure/Motivation/Hard_Samples.pdf}
\caption{The schematic diagram of hard samples.}
\label{hardsample}
\end{figure}

%下面写一下难样本挖掘方法分类，然后再举例。

 Hard sample mining has been widely used in computer vision and other fields \cite{shrivastava2016training,sun2019mvp,zhu2019distance}. %\textbf{Introduce the different classes of hard sample mining methods, then show examples}.
\citet{Suh_2019_CVPR} proposed a stochastic hard negative mining method that adopts class signatures to keep track of feature embedding online. The method improves image retrieval accuracy substantially. In \cite{Sun_2019_ICCV}, a weighted complete bipartite graph-based maximum-value perfect matching for mining the hard samples was proposed to relieve adverse optimization and sample imbalance problems. \citet{xue2019hard} applied a hard sample mining method based on an enhanced deep multiple instance learning to find the hard samples from unlabeled training data. 
However, few hard sample mining methods have been employed in the field of fault diagnosis for wind turbines, and the method proposed in this study fills the gap in the field.
% To the best of our knowledge, this paper is the first study to employ a hard sample mining framework in the field of wind turbine pitch system fault diagnosis.

\begin{figure*}[htb]
\centering
\includegraphics[width=0.95\linewidth]{Figure/methodology/proposed_method_20230803.pdf}
\caption{Overview of proposed HSMSCL}
\label{overview}
\end{figure*}


\begin{algorithm}[htb]%精简版代码
\caption{Hard Sample Mining Enabled Supervised Contrastive Learning}

\label{hsmsclalgorithm}
% \SetKwInput{KwIn}{Input}
% \SetKwInput{KwOut}{Output}
\SetKwProg{Function}{Function}{}{}
% \SetKwProg{For}{For}{:}{}
% \SetKwProg{While}{While}{:}{}
% \SetKw{Return}{return}

\KwIn{labeled train set $(x,y) \in D$, learning rate $\eta$, batch\_size $bs$, epochs $n$, temperature $\tau$, model $f$ and $g$ with parameters $\theta$ and $\theta_g$, number of HSM stages $num\_stages$}%, HSM threshold $p_1$ and $p_2$
\KwOut{trained model with parameters $\theta^{*}$}

Net, $\theta \leftarrow \text{initialize}(Net)$\;
$x \leftarrow \text{Normalize}(x)$\;

\Function{\text{HSM(bs, D, mis\_samples)}}
{
    hard\_batches $\leftarrow$ []\;
    \For{$b = 1$ \KwTo $length(D) // bs$}
    {
        batch $\leftarrow$ []\;
        batch $\leftarrow$ Initialization(batch, $D_i$, mis\_samples)\;
        left $\leftarrow$ bs - length(batch)\;
        \While{$left > 0$}
        {
            append(batch, random\_sample($D$))\;
            s $\leftarrow$ random\_sample(batch)\;
            append(batch, hard\_positive($s$))\;
            append(batch, hard\_negative($s$))\;
        }
        append(hard\_batches, batch)\;
    }
    \Return hard\_batches\;
}
\underline{\textbf{Stage 1:} \textit{SCL}}

\For{$stage = 1$ \KwTo $num\_stages$}
{
    hard\_batches $\leftarrow$ HSM(bs, D, [])\;
    \For{$epoch = 1$ \KwTo $n$}
    {
        \For{$(x_b,y_b) \in hard\_batches$}
        {
            $out_{b} \leftarrow f(\theta, x_{b})$\;
            $\mathcal{L} \leftarrow$ loss\_SCL$(y_b, out_b, \tau)$\;
            $\theta \leftarrow$ back\_propagation$(\mathcal{L}, \theta, \eta)$\;
        }
    }
}
\underline{\textbf{Stage 2:} \textit{MLP}}

mis\_samples $\leftarrow$ []\;

\For{$stage = 1$ \KwTo $num\_stages$}
{
    hard\_batches $\leftarrow$ HSM(bs, D, mis\_samples)\;
    \For{$epoch = 1$ \KwTo $n$}
    {
        \For{$(x_b,y_b) \in hard\_batches$}
        {
            $out_{b} \leftarrow g(\theta_g, x_{b})$\;
            $\mathcal{L} \leftarrow$ loss\_MLP$(y_b, out_b, \tau)$\;
            $\theta_g \leftarrow$ back\_propagation$(\mathcal{L}, \theta_g, \eta)$\;
        }
    }
    mis\_samples $\leftarrow$ misclassification(g, D)\;
}
\end{algorithm}

\section{Proposed method}
\label{proposedmethod}
\subsection{Overview}%画一张方法示意图
%将HSM方法与SCL方法结合进行故障诊断，有效解决了由于多分类不平衡问题导致的样本难分类问题。

The structure of the proposed method is shown in Fig.\ref{overview}. The proposed method consists of five parts: data preprocessing, data augmentation, hard sample mining, supervised contrastive learning, and multi-layer perceptron.

\begin{enumerate}%[label={\arabic*)}]%\alph* 表示小写英文字母,\arabic* 表示阿拉伯数字
\item Data preprocessing

\begin{comment}
    First, in order to eliminate the effect of the magnitudes of different variables, the original SCADA data need to be normalized.
\end{comment}
SCADA data is collected by hundreds of various types of sensors inside the wind turbine, therefore, the raw data needs to be normalized in order to eliminate the effects of different variable sizes. In this paper, we use Minimum-Maximum Normalization, which individually scales each data feature to [0,1] by a linear transformation of the original data. 
The formula of Minimum-Maximum Normalization is shown as follows:
\begin{equation}
\begin{aligned}
    x_i^{\text {norm }}=\frac{x_i-x_i^{\min }}{x_i^{\max }-x_i^{\min }}
\end{aligned}
\label{normalization}
\end{equation}
where $x_i$ is the $i$th feature of the input sample and $x_i^{\max}$ and $x_i^{\min}$ are the maximum and minimum values found in $x_i$, respectively.
% \textbf{cite the initial paper of minmax Normalization}

% \item Data augmentation(\textbf{consider delete it})

% Data augmentation methods can make the training data more diverse and thus improve the performance of the model\cite{zhao2020deep}. We have two data augmentation methods, namely RandomAddGaussian and RandomScale. RandomAddGaussian adds Gaussian noise randomly into the input data, while RandomScale multiplies the input data with a random factor. In the experiments, one of them is chosen randomly for the preprocessed data.

% RandomAddGaussian is the strategy that randomly adds Gaussian noise into the input data. The formulation of the RandomAddGaussian is as follows:
% \begin{equation}
%     x:=x+n
% \label{randomG}
% \end{equation}
% where $x$ is the input data, $n$ is the random Gaussian noise with the distribution $\mathcal{N}(0,0.01)$.

% RandomScale is the strategy that multiplies the input data with a random factor.  The formulation of the RandomScale is as follows:
% \begin{equation}
%     x:=\beta *x
% \label{randomS}
% \end{equation}
% where $x$ is the input data, $\beta$ is the random scaler with the distribution $\mathcal{N}(1,0.01)$.

\item Hard sample mining(HSM)

% Hard samples refer to samples that are difficult to be accurately classified by the model\cite{8588326}, including hard positive samples and hard negative samples.
In this paper, the proposed HSM method based on cosine similarity is used to select the hard samples. Specifically, when constructing a mini-batch, the proposed HSM method first randomly selects a sample from each class as initialization. Subsequently, either a random sample from the dataset is selected or a random query sample is chosen from the current mini-batch and its hard samples are determined to extend this mini-batch. In the latter case, the sample of a different class with the highest cosine similarity to the selected query sample is selected as its hard negative sample, while the sample of the same class with the lowest cosine similarity to the query sample is selected as a hard positive sample. We keep repeating the above process until the size of the mini-batch reaches the batch size. It is worth noting that the proposed HSM method is applied to construct the mini-batches of both the SCL and MLP stages. In the MLP stage, specifically, the misclassified samples for the current stage are used by HSM to initialize the mini-batches in the next stage. In this case, each time at initialization, a sample is randomly taken from misclassified samples if it is not empty, or from the dataset otherwise. 

% \item Supervised contrastive learning(SCL)
%讲一讲SCL的原理，SCL的优点。

\item Supervised Contrastive Learning(SCL)
%加梯度下降法训练的目标函数公式

% The proposed method trains the model by gradient descent method and the object function is shown below. 
With the powerful representation learning capability of SCL, the proposed method is capable of learning the rich representations in the data. In addition, the proposed method combines HSM with SCL by utilizing HSM to construct more challenging mini-batches, thus allowing SCL to learn more discriminative representations.

% Moreover, constructing more challenging mini-batches by HSM allows SCL to learn better representations.
 

\item Multi-layer perceptron(MLP)

The representations learned by SCL will be constructed by HSM into more challenging mini-batches for training MLP to classify the input data.

\end{enumerate}

\subsection{Hard sample mining enabled supervised contrastive learning(HSMSCL)}%具体介绍HSM与SCL的结合，绘制算法流程图
%HSM嵌入到SCL中学到了更好的表示，嵌入到MLP中训练出了更好的分类器,HSM_v2将每类数据数量都扩展到和数量最多的那一类一样。HSM方法之间难正样本计算方式为计算点与聚类中心的余弦相似度，难负样本计算方式为计算点与点之间的余弦相似度。

%20230706，对SCL阶段来说，HSM构建的是难分类数据样本；对MLP阶段来说，HSM构建的是难分类Embedding！！！！！！！！！！！！！！！！！！！！

The key to the proposed HSMSCL is to combine HSM as a unified framework with SCL and MLP, allowing SCL to learn more compact representations and train better MLP classifiers for more accurate diagnosis of hard-to-classify faults in wind turbine pitch systems. As shown in detail in Algorithm\ref{hsmsclalgorithm}, the two stages, SCL and MLP, each invoke HSM as a framework to construct more challenging mini-batches. 

By introducing hard samples through HSM, SCL enhances the learning of the model for sample similarity and difference, which in turn improves the model's performance in the fault diagnosis task. Through multiple rounds of training, the model gradually learns more discriminative features, which provides a better basis for feature representation in the subsequent MLP stage.

The main purpose of the MLP stage is to mine hard samples by using HSM and use these hard samples for training the MLP model. Hard samples introduce challenges in the training process and can motivate the model to better adapt to complex situations, thus improving the overall performance. Repeating the process of multiple rounds of training and updating the list of misclassified samples allows the MLP model to progressively optimize itself and enhance its accuracy and robustness to complex fault diagnosis tasks.

Overall, the proposed HSMSCL improves the performance of the fault diagnosis model for hard samples by utilizing the HSM framework based on the two-step approach of SCL and MLP.%对此段要做自己的修改

% \begin{algorithm}[htb]
% \caption{Hard sample mining enabled contrastive learning for wind turbine pitch system fault diagnosis}\label{hsmsclalgorithm}
% \KwIn{labeled train set $(x,y) \in D$, learning rate $\eta$, max epochs $n$, temperature $\tau$, model $Net$ with parameters $\theta$, number of HSM stages $hs$}
% \KwOut{trained model with parameters $\theta^{*}$}
% Net, $\theta$ $\leftarrow$ initialize(Net)\;  
% $x \leftarrow$ Normalize(x)$\;
% \\\For{ $stage$=$\{1...hs\}$} {
%     $hard\_batches$ \leftarrow $HSM(x)$\;
%     \For{$epoch$=$\{1...n\}$}
%     {   
%         \For{$(x_b,y_b)$ \in $hard\_batches$}
%         {
%         $out_{b}$ \leftarrow $Net(\theta, x_{b})$\;
%         \mathcal{L} \leftarrow$ loss\_SCL$(y_b, out_b, \theta, \tau)$\;
%         $grad_{\theta} \leftarrow$compute\_gradient$(\mathcal{L}, \theta)$\;
%         $\theta^{*} \leftarrow$ gradient\_descent$(Net,  \theta,grad_{  \theta},\eta)$\; 
%         $\theta$\leftarrow $\theta^{*}$ \;
%         }
%     }
% }
% \end{algorithm}

% \begin{algorithm}[htb]
% \caption{Hard sample mining enabled contrastive learning for wind turbine pitch system fault diagnosis}\label{hsmsclalgorithm}
% \KwIn{labeled train set $(x,y) \in D$, learning rate $\eta$, batch\_size $bs$, max epochs $n$, temperature $\tau$, model $Net$ with parameters $\theta$, number of HSM stages $num\_stages$, HSM threshold $p_1$ and $p_2$}
% \KwOut{trained model with parameters $\theta^{*}$}
% Net, $\theta$ $\leftarrow \text{initialize}(Net)$\\
% $x \leftarrow \text{Normalize}(x)$$\\
% \\\For{$stage=\{1...num\_stages\}$} {
%     hard\_batches $\leftarrow$ [] \\
%     \For{$b = \{1...length(D) // bs\}$}{
%         batch $\leftarrow$ [] \\
%         \For{$i = \{1...num\_classes\}$} {
%             $append(batch, random\_sample(D_i)$
%         }
%         $left \leftarrow bs - length(batch)$ \\
%         \While{$left > 0$} {
%             $p \leftarrow random [0, 1]$\\
%             \uIf{$p \in [0, p_1)$} {
%                 $append(batch, random\_sample(D))$
%             }
%             \uElseIf{$p \in [p_1, p_2)$} {
%                 s \leftarrow random\_sample(batch) \\
%                 $append(batch, hard\_positive(s))$
%             }
%             \Else {
%                 s \leftarrow random\_sample(batch) \\
%                 $append(batch, hard\_negative(s))$
%             }
%         }
%         $append(hard\_batches, batch)$ \\
%     }\
    
%     \For{$epoch=\{1...n\}$}
%     {   
%         \For{$(x_b,y_b) \in hard\_batches$}
%         {
%         $out_{b} \leftarrow$ $Net(\theta, x_{b})$\;
%         $\mathcal{L} \leftarrow$ loss\_SCL$(y_b, out_b, \theta, \tau)$\;
%         $grad_{\theta} \leftarrow$compute\_gradient$(\mathcal{L}, \theta)$\;
%         $\theta^{*} \leftarrow$ gradient\_descent$(Net,  \theta,grad_{  \theta},\eta)$\; 
%         $\theta\leftarrow \theta^{*}$ \;
%         }
%     }
% }
% \end{algorithm}


% \begin{algorithm}[htb]%精简版代码，我修改的
% \caption{Hard sample mining enabled contrastive learning for wind turbine pitch system fault diagnosis}\label{hsmsclalgorithm}
% \KwIn{labeled train set $(x,y) \in D$, learning rate $\eta$, batch\_size $bs$, epochs $n$, temperature $\tau$, model $f$ and $g$ with parameters $\theta$ and $\theta_g$,  number of HSM stages $num\_stages$, HSM threshold $p_1$ and $p_2$}
% \KwOut{trained model with parameters $\theta^{*}$}
% Net, $\theta$ $\leftarrow \text{initialize}(Net)$\\
% $x \leftarrow \text{Normalize}(x)$$\\

% \Function{HSM}{bs, D, mis\_samples} \\
%     \State hard\_batches $\leftarrow$ [] \\
%     \For{$b = \{1...length(D) // bs\}$}{
%          batch $\leftarrow$ [] \\
%         % \For{$i = \{1...num\_classes\}$} {
%         %      append(batch, initial\_sample($D_i$, anchors)
%         % }
%         batch $\leftarrow$ Initialization(batch, $D_i$, mis\_samples) \\
%         % left $\leftarrow$ bs - length(batch) \\
%         \While{$left > 0$} {
%             % $p \leftarrow random [0, 1]$\\
%             % \uIf{$p \in [0, p_1)$} {
%                  append(batch, random\_sample($D$)) \\
%             % }
%             % \uElseIf{$p \in [p_1, p_2)$} {
%                  s $\leftarrow$ random\_sample(batch) \\
%                  append(batch, hard\_positive($s$))\\
%             % }
%             % \Else {
%                 % s $\leftarrow$ random\_sample(batch) \\
%                  append(batch, hard\_negative($s$))\\
%             % }
%         }
%         append(hard\_batches, batch) \\
%     }
%   \State \Return hard\_batches
% \EndFunction

% \underline{\textbf{Stage 1:} \textit{CFL} } \\
% \\\For{$stage=\{1...num\_stages\}$} {
%     hard\_batches $\leftarrow$ HSM(bs, D, []) \\
%     \For{$epoch=\{1...n\}$}
%     {   
%         \For{$(x_b,y_b) \in hard\_batches$}
%         {
%         $out_{b} \leftarrow$ $f(\theta, x_{b})$\;
%         $\mathcal{L} \leftarrow$ loss\_SCL$(y_b, out_b, \tau)$\;
%         $\theta \leftarrow$ back\_propogation$(\mathcal{L}, \theta, \eta)$\;
%         }
%     }
% }\\

% \\\underline{\textbf{Stage 2:} \textit{MLP} } \\
% mis\_samples $\leftarrow$ []
% \\\For{$stage=\{1...num\_stages\}$} {
%     hard\_batches $\leftarrow$ HSM(bs, D, mis\_samples) \\
%     \For{$epoch=\{1...n\}$}
%     {   
%         \For{$(x_b,y_b) \in hard\_batches$}
%         {
%             $out_{b} \leftarrow$ $g(\theta_g, x_{b})$\;
%             $\mathcal{L} \leftarrow$ loss\_MLP$(y_b, out_b, \tau)$\;
%             $\theta_g \leftarrow$ back\_propogation$(\mathcal{L}, \theta_g, \eta)$\;
%         }
%     }
%     mis\_samples $\leftarrow$ misclassification(g, D)
% }\\
% \end{algorithm}


% \begin{algorithm}[htb]%完整版代码
% \caption{Hard sample mining enabled contrastive learning for wind turbine pitch system fault diagnosis}\label{hsmsclalgorithm}
% \KwIn{labeled train set $(x,y) \in D$, learning rate $\eta$, batch\_size $bs$, epochs $n$, temperature $\tau$, model $f$ and $g$ with parameters $\theta$ and $\theta_g$,  number of HSM stages $num\_stages$, HSM threshold $p_1$ and $p_2$}
% \KwOut{trained model with parameters $\theta^{*}$}
% Net, $\theta$ $\leftarrow \text{initialize}(Net)$\\
% $x \leftarrow \text{Normalize}(x)$$\\

% \Function{HSM}{bs, D, mis\_samples} \\
%     \State hard\_batches $\leftarrow$ [] \\
%     \For{$b = \{1...length(D) // bs\}$}{
%          batch $\leftarrow$ [] \\
%         % \For{$i = \{1...num\_classes\}$} {
%         %      append(batch, initial\_sample($D_i$, anchors)
%         % }
%         batch $\leftarrow$ Initialization(batch, $D_i$, mis\_samples) \\
%         left $\leftarrow$ bs - length(batch) \\
%         \While{$left > 0$} {
%             $p \leftarrow random [0, 1]$\\
%             \uIf{$p \in [0, p_1)$} {
%                  append(batch, random\_sample($D$)) \\
%             }
%             \uElseIf{$p \in [p_1, p_2)$} {
%                  s $\leftarrow$ random\_sample(batch) \\
%                  append(batch, hard\_positive($s$))
%             }
%             \Else {
%                 s $\leftarrow$ random\_sample(batch) \\
%                  append(batch, hard\_negative($s$))
%             }
%         }
%         append(hard\_batches, batch) \\
%     }
%   \State \Return hard\_batches
% \EndFunction

% \underline{\textbf{Stage 1:} \textit{SCL} } \\
% \\\For{$stage=\{1...num\_stages\}$} {
%     hard\_batches $\leftarrow$ HSM(bs, D, []) \\
%     \For{$epoch=\{1...n\}$}
%     {   
%         \For{$(x_b,y_b) \in hard\_batches$}
%         {
%         $out_{b} \leftarrow$ $f(\theta, x_{b})$\;
%         $\mathcal{L} \leftarrow$ loss\_SCL$(y_b, out_b, \tau)$\;
%         $\theta \leftarrow$ back\_propogation$(\mathcal{L}, \theta, \eta)$\;
%         }
%     }
% }\\

% \\\underline{\textbf{Stage 2:} \textit{MLP} } \\
% mis\_samples $\leftarrow$ []
% \\\For{$stage=\{1...num\_stages\}$} {
%     hard\_batches $\leftarrow$ HSM(bs, D, mis\_samples) \\
%     \For{$epoch=\{1...n\}$}
%     {   
%         \For{$(x_b,y_b) \in hard\_batches$}
%         {
%             $out_{b} \leftarrow$ $g(\theta_g, x_{b})$\;
%             $\mathcal{L} \leftarrow$ loss\_MLP$(y_b, out_b, \tau)$\;
%             $\theta_g \leftarrow$ back\_propogation$(\mathcal{L}, \theta_g, \eta)$\;
%         }
%     }
%     mis\_samples $\leftarrow$ misclassification(g, D)
% }\\
% \end{algorithm}



% \begin{table}[htb]%table*占双栏，此为双行版本表
%  \centering
%  \caption{26 parameters of SCADA data.}
%  \label{variable}
%  % \resizebox{\textwidth}{!}{ %\resizebox{宽度}{高度}{对象},此处"!"表示根据根据宽高比进行自适应缩放
%  \begin{tabular}{l|l}%l左，c中，r右
%   \toprule%表格顶部的粗线
%   No. & Parameter  \\
%   \midrule%表格中间的细分隔线
%   1  & Wind speed        \\
%   2  & Generator speed \\
%   3 & Active power          \\
%   4 & Wind direction \\
%   5 & Average wind direction within 25s \\
%   6 & Yaw position     \\
%   7 & Yaw speed           \\
%   8 & Angle of pitch 1 \\
%   9 & Angle of pitch 2 \\
%   10 & Angle of pitch 3 \\
%   11 & Speed of pitch 1 \\
%   12 & Speed of pitch 2 \\
%   13 & Speed of pitch 3 \\
%   14 & Temperature of pitch motor 1\\
%   15 &  Temperature of pitch motor 2\\
%   16 & Temperature of pitch motor 3\\
%   17 & Horizontal acceleration\\
%   18 & Vertical acceleration\\
%   19 & Environment temperature\\
%   20 & Internal temperature of nacelle\\
%   21 & Switching temperature of pitch 1\\
%   22 & Switching temperature of pitch 2\\
%   23 & Switching temperature of pitch 3\\
%   24 & DC current of pitch 1 switch charger\\
%   25 & DC current of pitch 2 switch charger\\
%   26 & DC current of pitch 3 switch charger\\
%   \bottomrule%表格底部的粗线
%  \end{tabular}
% \end{table}

\begin{table}[htb]
\centering
\caption{26 parameters of SCADA data.}
\label{variable}
\begin{tabular}{l|l}
\toprule%表格顶部的粗线
No. & Parameter                            \\ \hline
1   & Wind speed                           \\
2   & Generator speed                      \\
3   & Active power                         \\
4   & Wind direction                       \\
5   & Average wind direction within 25s    \\
6   & Yaw position                         \\
7   & Yaw speed                            \\
8   & Angle of pitch 1                     \\
9   & Angle of pitch 2                     \\
10  & Angle of pitch 3                     \\
11  & Speed of pitch 1                     \\
12  & Speed of pitch 2                     \\
13  & Speed of pitch 3                     \\
14  & Temperature of pitch motor 1         \\
15  & Temperature of pitch motor 2         \\
16  & Temperature of pitch motor 3         \\
17  & Horizontal acceleration              \\
18  & Vertical acceleration                \\
19  & Environment temperature              \\
20  & Internal temperature of nacelle      \\
21  & Switching temperature of pitch 1     \\
22  & Switching temperature of pitch 2     \\
23  & Switching temperature of pitch 3     \\
24  & DC current of pitch 1 switch charger \\
25  & DC current of pitch 2 switch charger \\
26  & DC current of pitch 3 switch charger \\ 
\bottomrule%表格底部的粗线
\end{tabular}
\end{table}


\section{Experiment}

\begin{figure}[htb]
\centering
\subfloat[]{\includegraphics[width=2.5in]{Figure/Data_description/Data_distribution_of_WT_23.pdf}%
\label{datadistribution23}}
\hfil
\subfloat[]{\includegraphics[width=2.5in]{Figure/Data_description/Data_distribution_of_WT_29.pdf}%
\label{datadistribution29}}
\caption{The data distribution of the two datasets. (a) WT 23 (b) WT 29}
\label{datadistribution}
\end{figure}

\begin{figure}[htb]
\centering
\subfloat[]{\includegraphics[width=2.5in]{Figure/Data_description/power3d_23.png}%
\label{power23}}
\hfil
\subfloat[]{\includegraphics[width=2.5in]{Figure/Data_description/power3d_29.png}%
\label{power29}}
\caption{The output power VS wind speed for WT 23 and WT 29. (a) WT 23 (b) WT 29}
\label{power}
\end{figure}

\subsection{Experimental Setup}
\subsubsection{Data Description}
\label{secdatadesc}
%数据集描述
To validate the effectiveness of the proposed method, two pitch system cog belt fracture datasets collected from two real-world wind turbines numbered 23 and 29 are used \cite{industrybigdata} (in the remainder of this paper, they are referred to as WT 23 and WT 29). A cog belt is a transmission device with a toothed structure on the working surface and is a key component in the electric pitch system of a wind turbine. During wind turbine operation, the cog belt may gradually break due to stress fatigue, leading to the failure of the wind turbine pitch system. There are five health conditions in the datasets: normal, slightly worn, low risk of fracture, high risk of fracture, and complete fracture. If the cog belt completely fractures, it can significantly affect the wind turbine's power production and increase maintenance costs. Therefore, it is crucial to accurately diagnose the health condition of the cog belt and replace it before it fractures completely to ensure the normal operation of the wind turbine. 

% Both datasets are recorded through the SCADA system and have 26 variables, as shown in Table \ref{variable}. The total number of data for WT 23 and WT 29 is 13,6709 and 215,012, respectively. Fig.\ref{datadistribution} shows the extreme imbalance in the duration of the different health conditions in the datasets. Fig.\ref{power} shows the output power VS wind speed for WT 23 and WT 29. It can be seen that the cog belt has different effects on the power generation efficiency of wind turbine when it is in different health conditions. During the experiment, we divided 70\% of the data into the training set and 30\% into the test set, where 20\%  of the training set is further divided as the validation set.

The original sensor data contains hundreds of variables recorded by a large number of sensors of the SCADA system inside the wind turbine. After selection by experts in the field, the remaining irrelevant parameters were removed and 26 parameters related to the cog belt fracture of the wind turbine were retained. The details of the parameters are shown in Table \ref{variable}. The total number of data for WT 23 and WT 29 is 13,6709 and 215,012, respectively. Fig.\ref{datadistribution} shows the extreme imbalance in the duration of the different health conditions in the datasets. Fig.\ref{power} shows the output power VS wind speed for WT 23 and WT 29. It can be seen that the cog belt has different effects on the power generation efficiency of wind turbine when it is in different health conditions. During the experiment, we divided 70\% of the data into the training set and 30\% into the test set, where 20\%  of the training set is further divided as the validation set.

% \textbf{Introduce the extreme class imbalance problem in the dataset. Add the power VS wind speed figure and explain the reason. Add the data preprocessing process.}



\subsubsection{Evaluation Metrics}%介绍准确率和macro Gmean
In this paper, In this paper, \text{Accuracy} and \text{macro G-mean} \cite{barandela2003strategies,ri2020g} are used as evaluation metrics. \text{G-mean} is a good indicator of the classifier's performance in imbalanced classification problems \cite{kubat1997addressing}. Due to the extreme class imbalance in the dataset used in this paper, the macro average of \text{G-mean} was used in order to fully represent the effectiveness of different methods in classifying each health condition. The evaluation metrics are defined as follows.


\begin{equation}
\text{Accuracy} = \frac{\mathrm{TP}+\mathrm{TN}}{\mathrm{TP+FP+FN+TN}}
\label{acc}
\end{equation}

\begin{equation}
\text{G-mean} = \sqrt{ \frac{\mathrm{TP}}{\mathrm{TP+FN}} \times \frac{\mathrm{TN}}{\mathrm{TN+FP}}}
\label{gmean}
\end{equation}

\begin{equation}
\operatorname{macro} \text{ G-mean}=\frac{1}{n} \sum_{i=1}^n \mathrm{G-mean_i}
\label{macrogmean}
\end{equation}
where \text{TP}, \text{TN}, \text{FP}, and \text{FN} refer to true positives, true negatives, false positives, and false negatives, respectively.

\subsubsection{Baseline}%介绍对比方法，将原来的架构做一些调整，在baseline中引入SMOTE或FL来解决样本不平衡问题；此外，SCL—MLP也作为baseline方法验证HSM架构的有效性

To verify the effectiveness of the proposed method, we selected 5 baselines, namely SCL-MLP (with the HSM framework removed from the proposed method), ResNet101, BiLSTM, and CNN. In order to solve the problem of extreme class imbalance in the datasets, Focal Loss (FL) \cite{lin2017focal} was used in ResNet101, BiLSTM, and CNN. A brief description of the baselines is given below.

\begin{enumerate}%[label={\arabic*)},topsep=0pt,itemsep=0pt,parsep=0pt,partopsep=0pt]
\item SCL-MLP is the method after removing the HSM framework from the proposed method.
\item ResNet101 is a residual network with a depth of 101 layers.
\item BiLSTM is a recurrent neural network consisting of a forward LSTM and a backward LSTM.
\item CNN used in this paper is a deep neural network consisting of 4 convolutional layers and 3 fully connected layers.
\item MLP used in this paper consists of three fully connected layers.
\end{enumerate}


\begin{figure*}[htbp]
\centering
\subfloat[]{\includegraphics[width=2in]{Figure/Experiment/TSNE_HSMstage4+SCL23_new.png}%
\label{TSNE23HSMSCL}}
\hfil
\subfloat[]{\includegraphics[width=2in]{Figure/Experiment/TSNE_SCL23.png}%
\label{TSNE23SCL}}
\hfil
\subfloat[]{\includegraphics[width=2in]{Figure/Experiment/TSNEMLPNum23.png}%
\label{TSNE23MLP}}
\hfil
\\
\subfloat[]{\includegraphics[width=2in]{Figure/Experiment/TSNEResNetNum23.png}%
\label{TSNE23Resnet}}
\hfil
\subfloat[]{\includegraphics[width=2in]{Figure/Experiment/TSNEBiLSTMNum23.png}%
\label{TSNE23LSTM}}
\hfil
\subfloat[]{\includegraphics[width=2in]{Figure/Experiment/TSNECNNNum23.png}%
\label{TSNE23CNN}}
\caption{T-SNE embedding for WT 23 (a) Proposed method (b) SCL+MLP (c) MLP+FL (d) ResNet101+FL (e) BiLSTM+FL (f) CNN+FL}
\label{TSNE23}
\end{figure*}

\begin{figure*}[htbp]
\centering
\subfloat[]{\includegraphics[width=2in]{Figure/Experiment/HSMSCL_num23_confusion_matrix_new.png}%
\label{confusion23HSMSCL}}
\hfil
\subfloat[]{\includegraphics[width=2in]{Figure/Experiment/confusion_matrix_SCL+MLP_Num23.png}%
\label{confusion23SCL}}
\hfil
\subfloat[]{\includegraphics[width=2in]{Figure/Experiment/confusion_matrix_MLP_Num23.png}%
\label{confusion23MLP}}
\hfil
\\
\subfloat[]{\includegraphics[width=2in]{Figure/Experiment/confusion_matrix_Resnet_Num23.png}%
\label{confusion23Resnet}}
\hfil
\subfloat[]{\includegraphics[width=2in]{Figure/Experiment/confusion_matrix_LSTM_Num23.png}%
\label{confusion23LSTM}}
\hfil
\subfloat[]{\includegraphics[width=2in]{Figure/Experiment/confusion_matrix_CNN_Num23.png}%
\label{confusion23CNN}}
\caption{Confusion matrix for WT 23 (a) Proposed method (b) SCL+MLP (c) MLP+FL (d) ResNet101+FL (e) BiLSTM+FL (f) CNN+FL}
\label{confusion23}
\end{figure*}

\subsection{HSMSCL}

\subsubsection{Experiment for WT 23}
 We compared the proposed HSMSCL and the selected 5 baselines on the dataset of WT 23. We repeated the experiment ten times with different random seeds and the reported results are based on the average. The results are shown in Table \ref{resultWT23}. The corresponding t-SNE embeddings and confusion matrix are shown in Fig.\ref{TSNE23} and Fig.\ref{confusion23}. Compared with the baselines, the $Accuracy$ of the proposed method is increased by 5.41\%, 0.30\%, 0.03\%, 1.17\%, 0.15\% and 0.42\%, and the \text{macro G-mean} of the proposed method is increased by 3.07\%, 7.02\%, 3.29\%, 6.70\%, 9.43\% and 2.27\%. The results in Table \ref{resultWT23} show that the proposed HSMSCL performs significantly better than baselines, especially in the value of \text{macro G-mean}. Fig.\ref{TSNE23} indicates that the proposed method is able to diagnose the fault of the pitch system of WT 23 excellently. The confusion matrix shown in Fig.\ref{confusion23} indicates that the proposed method can effectively solve the problem of difficulty in distinguishing multiple faults in the pitch system of WT 23.

% \begin{table}[htb]%20230716版本
%  \centering
%  \caption{Comparison of fault diagnosis methods performance for WT 23}
%  \label{resultWT23}
%  \begin{tabular}{lccc}%l左，c中，r右
%   \toprule%表格顶部的粗线
%   Algorithms & \text{Accuracy} & \text{macro G-mean} \\
%   \midrule%表格中间的细分隔线
%   \textbf{HSMCFL}  & \textbf{0.9979} & \textbf{0.9991}\\%v2,scl_stage5,lc_stage4(WindTurbine_cogbelt_Hard Sample Mining_test\HSM+SCL_HSM+LC\HSM+SCLstage4_HSM+LC)
%   CFL-MLP  & 0.9467 &  0.9693\\%tem0.1,last_best
%   MLP & 0.9949 & 0.9336 \\ 
%   MLP+FL & 0.9976 & 0.9673 \\%alpha_f=[0.2,0.05,3,1,1]
%   ResNet101+FL & 0.9864 & 0.9364\\
%   BiLSTM+FL & 0.9964 & 0.9130\\
%   CNN+FL & 0.9937 &  0.9769\\
%   \bottomrule%表格底部的粗线
%  \end{tabular}
% \end{table}

\begin{table}[htb]%20230717版本
 \centering
 \caption{Comparison of fault diagnosis methods performance for WT 23}
 \label{resultWT23}
 \begin{tabular}{lccc}%l左，c中，r右
  \toprule%表格顶部的粗线
  Algorithms & \text{Accuracy} & \text{macro G-mean} \\
  \midrule%表格中间的细分隔线
  \textbf{HSMSCL}  & \textbf{0.9979} & \textbf{0.9991}\\%v2,scl_stage5,lc_stage4(WindTurbine_cogbelt_Hard Sample Mining_test\HSM+SCL_HSM+LC\HSM+SCLstage4_HSM+LC)
  SCL-MLP  & 0.9467 &  0.9693\\%tem0.1,last_best
  % MLP & 0.9949 & 0.9336 \\ 
  MLP+FL & 0.9759 & 0.8982 \\%alpha_f=[0.2,0.05,3,1,1]
  ResNet101+FL & 0.9728 & 0.9634\\
  BiLSTM+FL & 0.9129 & 0.9351\\
  CNN+FL & 0.9501 &  0.9509\\
  \bottomrule%表格底部的粗线
 \end{tabular}
\end{table}


\begin{figure*}[htbp]
\centering
\subfloat[]{\includegraphics[width=2in]{Figure/Experiment/TSNE_HSMstage4+SCL29_new.png}%
\label{TSNE29HSMSCL}}
\hfil
\subfloat[]{\includegraphics[width=2in]{Figure/Experiment/TSNE_SCL29.png}%
\label{TSNE29SCL}}
\hfil
\subfloat[]{\includegraphics[width=2in]{Figure/Experiment/TSNEMLPNum29.png}%
\label{TSNE29MLP}}
\hfil
\\
\subfloat[]{\includegraphics[width=2in]{Figure/Experiment/TSNEResNetNum29.png}%
\label{TSNE29Resnet}}
\hfil
\subfloat[]{\includegraphics[width=2in]{Figure/Experiment/TSNEBiLSTMNum29.png}%
\label{TSNE29LSTM}}
\hfil
\subfloat[]{\includegraphics[width=2in]{Figure/Experiment/TSNECNNNum29.png}%
\label{TSNE29CNN}}
\caption{T-SNE embedding for WT 29 (a) Proposed method (b) SCL+MLP (c) MLP+FL (d) ResNet101+FL (e) BiLSTM+FL (f) CNN+FL}
\label{TSNE29}
\end{figure*}

\begin{figure*}[htbp]
\centering
\subfloat[]{\includegraphics[width=2in]{Figure/Experiment/HSMSCL_num29_confusion_matrix_new.png}%
\label{confusion29HSMSCL}}
\hfil
\subfloat[]{\includegraphics[width=2in]{Figure/Experiment/confusion_matrix_SCL+MLP_Num29.png}%
\label{confusion29SCL}}
\hfil
\subfloat[]{\includegraphics[width=2in]{Figure/Experiment/confusion_matrix_MLP_Num29.png}%
\label{confusion29MLP}}
\hfil
\\
\subfloat[]{\includegraphics[width=2in]{Figure/Experiment/confusion_matrix_Resnet_Num29.png}%
\label{confusion29Resnet}}
\hfil
\subfloat[]{\includegraphics[width=2in]{Figure/Experiment/confusion_matrix_LSTM_Num29.png}%
\label{confusion29LSTM}}
\hfil
\subfloat[]{\includegraphics[width=2in]{Figure/Experiment/confusion_matrix_CNN_Num29.png}%
\label{confusion29CNN}}
\caption{Confusion matrix for WT 29 (a) Proposed method (b) SCL+MLP (c) MLP+FL (d) ResNet101+FL (e) BiLSTM+FL (f) CNN+FL}
\label{confusion29}
\end{figure*}

\subsubsection{Experiment for WT 29}

To verify the generalization performance of the proposed method, we also compared it with baselines on WT 29. We repeated the experiment ten times with different random seeds and the reported results are based on the average. The results are shown in Table \ref{resultWT29}. The corresponding t-SNE embeddings and confusion matrix are shown in Fig.\ref{TSNE29} and Fig.\ref{confusion29}.  Compared with the baselines, the \text{Accuracy} of the proposed method is increased by 9.46\%, 0.45\%, 4.97\%, 2.07\%, 1.68\% and 1.09\%, and the \text{macro G-mean} of the proposed method is increased by 4.39\%, 7.85\%, 3.01\%, 14.02\%, 2.31\% and 6.02\%. The results in Table \ref{resultWT29}, Fig. \ref{TSNE29} and Fig. \ref{confusion29} once again verify that the proposed method has excellent hard-sample classification capability, which means that the proposed method is applicable to the hard-to-classify problem present in the wind turbine pitch system.

% \begin{table}[htb]%20230715结果
%  \centering
%  \caption{Comparison of fault diagnosis methods performance for WT 29}
%  \label{resultWT29}
%  \begin{tabular}{lccc}%l左，c中，r右
%   \toprule%表格顶部的粗线
%   Algorithms & \text{Accuracy} & \text{macro G-mean} \\
%   \midrule%表格中间的细分隔线
%   \textbf{HSMCFL}  & \textbf{0.9940} & \textbf{0.9971}\\%scl_stage4,lc_stage4(Num29 experiment\HSM+SCL_HSM+MLPorLC_2nd)
%   CFL-MLP  & 0.9081 &  0.9552\\%last_tem0.2,final_best_tem0.2
%   MLP & 0.9895 & 0.9245 \\ 
%   MLP+FL & 0.9469 & 0.9680 \\ %alpha_f=[1,1,2,10,2]
%   ResNet101+FL & 0.9738 & 0.8745\\
%   BiLSTM+FL & 0.9776 & 0.9746\\
%   CNN+FL & 0.9833 &  0.9405\\
%   \bottomrule%表格底部的粗线
%  \end{tabular}
% \end{table}

\begin{table}[htb]%20230718结果
 \centering
 \caption{Comparison of fault diagnosis methods performance for WT 29}
 \label{resultWT29}
 \begin{tabular}{lccc}%l左，c中，r右
  \toprule%表格顶部的粗线
  Algorithms & \text{Accuracy} & \text{macro G-mean} \\
  \midrule%表格中间的细分隔线
  \textbf{HSMSCL}  & \textbf{0.9940} & \textbf{0.9971}\\%scl_stage4,lc_stage4(Num29 experiment\HSM+SCL_HSM+MLPorLC_2nd)
  SCL-MLP  & 0.9081 &  0.9552\\%last_tem0.2,final_best_tem0.2
  % MLP & 0.9895 & 0.9245 \\ 
  MLP+FL & 0.9945 & 0.9062 \\ 
  ResNet101+FL & 0.9719 & 0.9495\\
  BiLSTM+FL & 0.9867 & 0.9840\\
  CNN+FL & 0.9605 &  0.9180\\
  \bottomrule%表格底部的粗线
 \end{tabular}
\end{table}

\subsection{Impact of hyperparameter temperature $\tau$}
%讲temperature是对比学习中一个重要的参数，起到了什么作用，我们选取的temperature范围是[0,1]，经过实验，所提出的方法在temperature取值不同时其表现都很出色，在temperature为多少时表现最好。如图所示，所提出方法的故障诊断性能对于不同的temperature具有出色的鲁棒性。

The hyperparameter temperature is commonly used in contrastive learning \cite{wang2021understanding}, we have performed experiments on the proposed method by selecting different temperature values including 0.1, 0.2, 0.3, 0.4, 0.5, 0.6, 0.7, 0.8, 0.9 and 1.0. As shown in Fig. \ref{sensitivitytem}, the proposed fault diagnosis method's performance is excellent at different temperature values, and the best performance is achieved when the temperature is 0.8. Therefore, the fault diagnosis performance of the proposed method is robust to different values of temperature $\tau$.

\begin{figure}[htb]
\centering
\includegraphics[width=0.95\linewidth]{Figure/Experiment/Sensitivity_analysis_temperature.pdf}
\caption{macro G-mean as a function of the temperature parameter $\tau$}
\label{sensitivitytem}
\end{figure}

% \subsection{Impact of HSM threshold $p_1$ and $p_2$}
% \subsection{Influence of imbalance factor}
%五种类别的话怎么调节其不平衡比例呢？

\begin{table}[htb]
 \centering
 \caption{Ablation analysis of HSM for WT 23}
 \label{ablation23}
 \begin{tabular}{lccc}%l左，c中，r右
  \toprule%表格顶部的粗线
  Algorithms &  \text{macro G-mean} \\
  \midrule%表格中间的细分隔线
  \textbf{HSMSCL}   & \textbf{0.9991}\\%scl_stage5,lc_stage4
  HSM+SCL-MLP  &   0.9887\\%v2result20230622,scl_stage3,best_best
  SCL-HSM+MLP &  0.9808\\%v1result,lc_stage=1,best_best
  \bottomrule%表格底部的粗线
 \end{tabular}
\end{table}

\begin{table}[htb]
 \centering
 \caption{Ablation analysis of HSM for WT 29}
 \label{ablation29}
 \begin{tabular}{lccc}%l左，c中，r右
  \toprule%表格顶部的粗线
  Algorithms &  \text{macro G-mean} \\
  \midrule%表格中间的细分隔线
  \textbf{HSMSCL}  &  \textbf{0.9971}\\%scl_stage4,lc_stage4
  HSM+SCL-MLP  &   0.9682\\%scl_stage1,last_last
  SCL-HSM+MLP &  0.9595\\%v1result,lc_stage=1,best_best
  \bottomrule%表格底部的粗线
 \end{tabular}
\end{table}

\subsection{Ablation analysis of HSM}



To verify the effectiveness of the HSM framework for the two stages (SCL and MLP) in the proposed method, 
we compare the proposed method with the following two methods and evaluate their performance by \text{macro G-mean}.

HSM+SCL-MLP: The HSM framework is removed from MLP in the proposed method.

SCL-HSM+MLP: The HSM framework is removed from SCL in the proposed method.

The results of the ablation analysis are shown in Table \ref{ablation23} and Table \ref{ablation29}. From these two tables, it can be seen that \text{macro G-mean} of the proposed method is significantly higher than that of the two compared methods. Taking the experimental results on WT 23 shown in Table \ref{ablation23} as an example, the \text{macro G-mean} of the proposed method is 1.05\% and 1.87\% higher than that of HSM+SCL-MLP and SCL-HSM+MLP, respectively. The above results demonstrate the necessity of the HSM framework for the proposed method. 



\section{Conclusion}

In this paper, hard sample mining enabled unified contrastive feature learning is proposed to diagnose the fault in wind turbine pitch systems. The existence of hard samples between different health conditions is a practical problem in wind turbine pitch systems. The proposed method solves this problem by combining a hard sample mining framework with a two-step approach consisting of supervised contrastive learning and multi-layer perceptron. 
\begin{comment}
    Notably, this is the first time that a hard sample mining-based method is used for fault diagnosis of a wind turbine pitch system. 
\end{comment}
Notably, few previous works have focused on the problem of hard samples in the field of wind turbine fault diagnosis, with this paper being a pioneer in applying hard sample mining techniques to this field. The effectiveness of the proposed method is validated on two real wind turbine pitch system failure datasets. As demonstrated in the experiments, the proposed HSMSCL outperforms the baselines. This means that the proposed method introduces a novel and effective solution for wind turbine pitch system fault diagnosis.

%% References with bibTeX database:
% \newpage
% \bibliographystyle{model1-num-names}
\bibliographystyle{unsrtnat}%{}中plain的意思是用作者的姓名排序，unsrt按照引用顺序来排序
\bibliography{ref.bib}

\newpage

 
% \vspace{11pt}

% \begin{IEEEbiography}[{\includegraphics[width=1in,height=1.25in,clip,keepaspectratio]{Figure/Author/zixuan.jpg}}]
% {Zixuan Wang} is currently a doctoral student at College of Biomedical Engineering and Instrument Science in Zhejiang University, Hangzhou, China. He received his Bachelor degree at College of Automation from Nanjing University of Science and Technology in 2020. His research interests are process monitoring and prognostic and health management. \\
% \end{IEEEbiography}

% \vspace{-20pt}
% \begin{IEEEbiography}[{\includegraphics[width=1in,height=1.25in,clip,keepaspectratio{Figure/Author/boqin.jpg}}]
% {Bo Qin} is currently a Master Degree Candidate at ZJU-UIUC Institute in Zhejiang University, Hangzhou, China. He received his Bachelor degree of Information Science and Electrical Engineering from Zhejiang University in 2020. His research interest is artificial intelligence and fault diagnosis. \\
% \end{IEEEbiography}

% \vspace{-20pt}
% \includegraphics[width=1in,height=1.25in,clip,keepaspectratio{Figure/Author/peng.jpg.jpg}
% \vspace{-20pt}
% % \end{wrapfigure}
% % \par
% \noindent \textbf{Peng Peng} is currently a doctoral student at the National Engineering Research Centre of Computer Integrated Manufacturing System (CIMS-ERC) in Tsinghua University, Beijing, China.  He received his Bachelor degree at Department of Automation from Northeastern University in 2016. His research interests are process monitoring and prognostic and health management. \\

% \begin{wrapfigure}{l}{25mm} 
% \vspace{-20pt}
% \includegraphics[width=1in,height=1.25in,clip,keepaspectratio]{Figure/Author/hongwei.jpg}
% \vspace{-20pt}
% \end{wrapfigure}
% \par
% \noindent \textbf{Hongwei Wang} received the B.S. degree in information technology and instrumentation from Zhejiang University, China, in 2004, the M.S. degree in control science and engineering from Tsinghua University, China, in 2007, and the Ph.D. degree in design knowledge retrieval from the University of Cambridge, in 2012. From 2011 to 2018, he was a Lecturer and then, a Senior Lecturer in engineering design with the University of Portsmouth. He is currently a Tenured Professor with Zhejiang University and the University of Illinois at Urbana–Champaign Joint Institute, Zhejiang University. His research interests include knowledge engineering, industrial knowledge graph, intelligent and collaborative systems, and data-driven fault diagnosis. His research in these areas has been published in over 120 peer-reviewed papers in well-established journals and international conferences.  \\








\vfill

\end{document}


